\apendice{Especificación de Requisitos}

\section{Introducción}
En esta sección se documenta el proceso de ingeniería de requisitos llevado a
cabo durante las fases iniciales del proyecto. A partir del análisis del
problema y de las necesidades del dominio analítico, se han establecido los
objetivos principales del sistema y se ha redactado un catálogo formal
detallando las funcionalidades esperadas (\acrfullpl{rf}) y las restricciones
técnicas impuestas a la arquitectura (\acrfullpl{rnf}).

Finalmente, se modela el comportamiento del sistema desde la perspectiva de sus
usuarios mediante la especificación de \acrfullpl{cu}.

\subsection{Términos del Dominio}
Para facilitar la comprensión del catálogo de requisitos y de los flujos de
interacción del sistema, se definen a continuación los conceptos clave
relativos al dominio de datos analizado:

\begin{description}
	\item[\textbf{Trabajo de Fin de Grado (TFG):}] Proyecto académico y práctico obligatorio que los estudiantes de último curso de Ingeniería Informática desarrollan y defienden para demostrar las competencias y conocimientos adquiridos a lo largo de su titulación.
	\item[\textbf{Readme (Fichero README):}] Archivo de documentación (usualmente en formato Markdown o texto plano) ubicado en la raíz del repositorio de GitHub que proporciona información introductoria del proyecto, instrucciones de configuración, despliegue y uso de la aplicación desarrollada.
	\item[\textbf{Abstract (Resumen académico):}] Síntesis concisa del trabajo (habitualmente redactada tanto en español como en inglés) que resume de forma compacta los objetivos, la metodología empleada y las conclusiones principales del \Acrshort{tfg}.
	\item[\textbf{Issue (Incidencia/Tarea):}] Herramienta integrada en el repositorio de GitHub para registrar, planificar y realizar el seguimiento de tareas pendientes, errores de programación, debates de diseño o hitos del desarrollo. En este contexto, representan el histórico de actividad y comunicación del proyecto.
	\item[\textbf{Contenido / Memoria del TFG:}] Texto íntegro contenido en los ficheros fuente en \LaTeX{} redactados por el alumno, que constituye el documento analítico y descriptivo formal de su trabajo.
\end{description}

\section{Objetivos generales}
El principal propósito de \nombreProyecto{} es proporcionar una herramienta que
facilite el análisis temático de los \acrshort{tfg} del grado de Ingeniería
Informática de la Universidad de Burgos disponibles en repositorios públicos de
GtHub. Para lograrlo, se han definido los siguientes objetivos generales:

\begin{itemize}
	\item \textbf{Extracción de datos:} Extraer el contenido de las memorias de \acrshort{tfg}, \textit{readmes}, \textit{issues} y \textit{abstracts} de repositorios públicos de GitHub.
	\item \textbf{Implementar procesamiento \acrshort{nlp}:} Normalizar los textos extraídos y prepararlos para el modelado temático.
	\item \textbf{Entrenar y evaluar algoritmos de modelado de temas:} Entrenar modelos como \acrshort{lda}, BERTopic, Top2Vec y FASTopic para la extracción de temas.
	\item \textbf{Desarrollar una interfaz web:} Permitir a los investigadores explorar de forma interactiva la distribución de temas y comparar empíricamente el rendimiento de cada modelo.
\end{itemize}

\section{Catálogo de requisitos}
El catálogo de requisitos desglosa los objetivos en unidades atómicas y
verificables, dividiéndolas según su naturaleza.

\subsection{\acrfullpl{rf}}
Definen las funcionalidades y acciones concretas que el sistema debe ser capaz
de ejecutar.

\begin{itemize}
	\item \textbf{RF-1. Extracción de datos:} El sistema debe descargar la información de repositorios a partir de un listado de orígenes parametrizado, comunicándose con la \acrshort{api} de GitHub.
	\item \textbf{RF-2. Preprocesamiento de texto:} El sistema debe aplicar técnicas de procesamiento de lenguaje natural (tokenización, lematización, borrado de \textit{stopwords}) para normalizar el texto.
	\item \textbf{RF-3. Entrenamiento de modelos:} El sistema debe ser capaz de instanciar y entrenar múltiples algoritmos de modelado de temas sobre el texto preprocesado.
	\item \textbf{RF-4. Optimización de hiperparámetros:} El sistema debe integrar mecanismos para evaluar y afinar los hiperparámetros de los modelos en busca de la máxima coherencia semántica.
	\item \textbf{RF-5. Cálculo de métricas:} El sistema debe generar métricas estandarizadas para evaluar la coherencia y diversidad de los temas.
	\item \textbf{RF-6. Visualización de temas:} La interfaz de usuario debe mostrar gráficas interactivas que representen la importancia de las palabras y de los temas.
	\item \textbf{RF-7. Comparativa de modelos:} La interfaz debe disponer de un panel que cruce y ordene el rendimiento métrico de los diferentes algoritmos.
	\item \textbf{RF-8. Análisis exploratorio académico:} La interfaz de usuario debe permitir filtrar los proyectos (por tutor, rango de años y calificaciones) y visualizar gráficos sobre la distribución y evolución de temáticas por curso académico.
\end{itemize}

\subsection{\acrfullpl{rnf}}
Restricciones de diseño, rendimiento y arquitectura que garantizan la calidad y
escalabilidad del producto final.

\begin{itemize}
	\item \textbf{RNF-1. Persistencia optimizada:} Dado el alto volumen de datos textuales, el almacenamiento local entre los módulos del sistema debe realizarse exclusivamente en formato Apache Parquet.
	\item \textbf{RNF-2. Desacoplamiento estructural:} La lógica de negocio del sistema debe aislarse de los detalles técnicos y librerías externas siguiendo el patrón de Arquitectura Hexagonal.
	\item \textbf{RNF-3. Contenerización:} Para evitar conflictos de dependencias, la ejecución de los diferentes módulos debe estar soportada por contenedores Docker independientes.
	\item \textbf{RNF-4. Tolerancia a fallos de red:} El módulo de ingesta debe gestionar de forma silenciosa la posible pérdida de conexión o la caída de la \acrshort{api} remota mediante reintentos y persistencia parcial.
\end{itemize}

\section{Especificación de \acrshortpl{cu}}
La especificación de \acrfullpl{cu} modela las interacciones entre los usuarios (o sistemas externos) y la aplicación para lograr objetivos específicos.

\subsection{Catálogo de Actores}
Para acotar el alcance de las interacciones, se definen a continuación los
roles o sistemas externos que participan en el sistema analítico:

\begin{description}
	\item[\textbf{Administrador / Desarrollador (Humano, Primario Backend):}] Perfil con conocimientos técnicos encargado de la configuración del sistema, la puesta en marcha de los entornos y la ejecución desatendida del pipeline de datos en sus etapas de ingesta y procesamiento desde la CLI.
	\item[\textbf{Usuario Académico (Humano, Primario Frontend):}] Docente o estudiante que busca respuestas funcionales en el corpus analizado. Su objetivo es identificar qué tecnologías se investigan más, cómo evolucionan las temáticas a lo largo de los cursos y qué tendencias existen en los \ACRshortpl{tfg} de Ingeniería Informática, todo ello a través de filtros y distribuciones temporales sencillas y sin necesidad de conocer en detalle el modelado de temas.
	\item[\textbf{Usuario Técnico / Investigador IA (Humano, Primario Frontend):}] Experto en \acrshort{nlp} que evalúa la calidad científica y estadística de los resultados. Utiliza la interfaz web técnica para analizar proyecciones, comprobar curvas de optimización de hiperparámetros y comparar cuantitativamente las métricas de coherencia entre los algoritmos.
	\item[\textbf{API de GitHub (Sistema Externo, Secundario):}] Interfaz de programación provista por GitHub que actúa como proveedor de datos externo. Es consumida por el pipeline para descargar de forma automatizada el contenido de las memorias de \acrshort{tfg}, \textit{readmes}, \textit{issues} y \textit{abstracts} del repositorio semilla.
\end{description}

\subsection{Diagrama de Casos de Uso}
El panorama general de interacciones queda plasmado en el diagrama de
\acrshortpl{cu} adjunto (Figura~\ref{fig:diagramas/casos_uso.png}).

\imagen{diagramas/casos_uso.png}{Diagrama de \acrshortpl{cu} del sistema \nombreProyecto}

\subsection{Fichas Tabulares de Casos de Uso}
A continuación, se documentan de forma tabular los flujos detallados asociados
a cada uno de estos \acrshortpl{cu}.

% Caso de Uso 1
\tablaSmallSinColores{CU-1: Ejecutar pipeline de ingesta.}
{p{0.21\textwidth} p{0.71\textwidth}}
{cu_1}
{
	\textbf{CU-1}                 & \textbf{Ejecutar pipeline de ingesta}                                                                                \\
}
{
	\textbf{Versión}              & 1.0                                                                                                                  \\
	\textbf{Autor}                & Alumno                                                                                                               \\
	\textbf{Actores}              & \textbf{Principal:} Administrador / Desarrollador. \newline \textbf{Secundario:} API de GitHub (Sistema externo).    \\
	\textbf{Requisitos asociados} & RF-1, RNF-1, RNF-4                                                                                                   \\
	\textbf{Descripción}          & El administrador inicia la recolección de los datos desde la nube.                                                   \\
	\textbf{Precondición}         & Disponer de conexión a Internet, de un archivo semilla de repositorios y de un token de GitHub válido.               \\
	\textbf{Acciones}             &
	\begin{enumerate}
		\def\labelenumi{\arabic{enumi}.}
		\tightlist
		\item El administrador lanza la orden de ejecución.
		\item El sistema lee el fichero semilla.
		\item El sistema consulta la \acrshort{api} y descarga el contenido de las memorias
		      de \acrshort{tfg}, \textit{readmes}, \textit{issues} y \textit{abstracts}.
		\item El sistema formatea los datos y los persiste.
	\end{enumerate}                                                                     \\
	\textbf{Postcondición}        & Se han generado los archivos con los datos extraídos.                                                    \\
	\textbf{Excepciones}          & Fallo de red, token expirado o repositorios eliminados (se registran en los \textit{logs} sin detener la ejecución). \\
	\textbf{Importancia}          & Alta                                                                                                                 \\
}

% Caso de Uso 2
\tablaSmallSinColores{CU-2: Entrenar modelos de temas.}
{p{0.21\textwidth} p{0.71\textwidth}}
{cu_2}
{
	\textbf{CU-2}                 & \textbf{Entrenar modelos de temas}                                                                            \\
}
{
	\textbf{Versión}              & 1.0                                                                                                             \\
	\textbf{Autor}                & Alumno                                                                                                          \\
	\textbf{Actores}              & \textbf{Principal:} Administrador / Desarrollador.                                                              \\
	\textbf{Requisitos asociados} & RF-2, RF-3, RF-4, RF-5                                                                                          \\
	\textbf{Descripción}          & El administrador ejecuta el proceso de análisis de los datos extraídos.          \\
	\textbf{Precondición}         & Deben existir archivos de ingesta generados por el CU-1.                                                        \\
	\textbf{Acciones}             &
	\begin{enumerate}
		\def\labelenumi{\arabic{enumi}.}
		\tightlist
		\item El administrador lanza el módulo de \textit{processing}.
		\item El sistema carga en memoria los textos crudos y aplica los filtros de limpieza
		      \acrshort{nlp}.
		\item El sistema entrena iterativamente los diferentes algoritmos.
		\item El sistema optimiza automáticamente los hiperparámetros de los modelos.
		\item Se calculan las métricas para la vista final.
	\end{enumerate}                                                             \\
	\textbf{Postcondición}        & Se han generado los ficheros con los resultados analíticos correspondiente a cada modelo. \\
	\textbf{Excepciones}          & Insuficiente memoria RAM para cargar un corpus completo (se aborta el hilo).                                    \\
	\textbf{Importancia}          & Alta                                                                                                            \\
}

% Caso de Uso 3
\tablaSmallSinColores{CU-3: Explorar temas de un modelo.}
{p{0.21\textwidth} p{0.71\textwidth}}
{cu_3}
{
	\textbf{CU-3}                 & \textbf{Explorar temas de un modelo}                                                                    \\
}
{
	\textbf{Versión}              & 1.0                                                                                                       \\
	\textbf{Autor}                & Alumno                                                                                                    \\
	\textbf{Actores}              & \textbf{Principal:} Usuario Técnico / Investigador IA.                                                    \\
	\textbf{Requisitos asociados} & RF-6                                                                                                      \\
	\textbf{Descripción}          & El usuario técnico visualiza las temáticas descubiertas por un algoritmo en particular de forma interactiva. \\
	\textbf{Precondición}         & El sistema web está arrancado y los modelos han sido entrenados. Deben existir archivos de ingesta generados por el CU-1. Deben existir archivos generados por el CU-2.                                          \\
	\textbf{Acciones}             &
	\begin{enumerate}
		\def\labelenumi{\arabic{enumi}.}
		\tightlist
		\item El usuario accede a la pestaña de análisis de modelo en el \textit{frontend}.
		\item Selecciona mediante un desplegable la fuente de texto y el modelo objetivo.
		\item El sistema carga bajo demanda la información del modelo.
		\item El usuario interactúa con los gráficos.
	\end{enumerate}                                                        \\
	\textbf{Postcondición}        & El usuario adquiere conclusiones sobre el corpus.                                                         \\
	\textbf{Excepciones}          & Datos no disponibles si el modelo no fue entrenado (el \textit{frontend} muestra advertencia).            \\
	\textbf{Importancia}          & Alta                                                                                                      \\
}

% Caso de Uso 4
\tablaSmallSinColores{CU-4: Comparar rendimiento de modelos.}
{p{0.21\textwidth} p{0.71\textwidth}}
{cu_4}
{
	\textbf{CU-4}                 & \textbf{Comparar rendimiento de modelos}                                                                              \\
}
{
	\textbf{Versión}              & 1.0                                                                                                                   \\
	\textbf{Autor}                & Alumno                                                                                                                \\
	\textbf{Actores}              & \textbf{Principal:} Usuario Técnico / Investigador IA.                                                                \\
	\textbf{Requisitos asociados} & RF-7                                                                                                                  \\
	\textbf{Descripción}          & El usuario técnico coteja las métricas de distintos algoritmos para evaluar cuál se comporta mejor ante el conjunto de datos. \\
	\textbf{Precondición}         & Al menos dos modelos diferentes deben haber sido entrenados para la misma fuente de datos.                            \\
	\textbf{Acciones}             &
	\begin{enumerate}
		\def\labelenumi{\arabic{enumi}.}
		\tightlist
		\item El usuario accede a la pestaña de comparativa.
		\item El sistema agrega la puntuación de coherencia de todos los modelos disponibles.
		\item Se renderiza una tabla de clasificación o \textit{ranking} que consolida el
		      rendimiento.
		\item El usuario puede visualizar gráficos comparativos.
	\end{enumerate}                                                                  \\
	\textbf{Postcondición}        & Ninguna alteración del estado del sistema.                                                                            \\
	\textbf{Excepciones}          & Ninguna.                                                                                                              \\
	\textbf{Importancia}          & Media                                                                                                                 \\
}

% Caso de Uso 5
\tablaSmallSinColores{CU-5: Consultar tendencias y metadatos académicos.}
{p{0.21\textwidth} p{0.71\textwidth}}
{cu_5}
{
	\textbf{CU-5}                 & \textbf{Consultar tendencias y metadatos académicos}                                                                 \\
}
{
	\textbf{Versión}              & 1.0                                                                                                                  \\
	\textbf{Autor}                & Alumno                                                                                                               \\
	\textbf{Actores}              & \textbf{Principal:} Usuario Académico.                                                                               \\
	\textbf{Requisitos asociados} & RF-8                                                                                                                 \\
	\textbf{Descripción}          & El usuario académico filtra los proyectos y analiza visualmente la evolución de temas y tecnologías por curso.       \\
	\textbf{Precondición}         & La interfaz web está iniciada y dispone de los datos de los proyectos cargados.                                  \\
	\textbf{Acciones}             &
	\begin{enumerate}
		\def\labelenumi{\arabic{enumi}.}
		\tightlist
		\item El usuario accede a la pestaña de análisis académico en el \textit{frontend}.
		\item Aplica filtros por rango de años, notas o nombre del tutor si lo desea.
		\item El sistema actualiza el listado y muestra las fichas de los proyectos
		      correspondientes.
		\item El usuario visualiza la distribución de temas y evolución de tecnologías en los
		      gráficos interactivos de la sección.
	\end{enumerate}                                                                                                                                     \\
	\textbf{Postcondición}        & Se presenta la información filtrada y los gráficos de evolución correspondientes.                                    \\
	\textbf{Excepciones}          & Ninguna.                                                                                                             \\
	\textbf{Importancia}          & Alta                                                                                                                 \\
}

\section{Matriz de trazabilidad}
Para asegurar que todos los \acrshortpl{rf} identificados están cubiertos por
las interacciones de los usuarios con el sistema, se presenta a continuación la
matriz de trazabilidad entre los \acrshortpl{cu} y los \acrshortpl{rf}.

\tablaSmall{Matriz de trazabilidad entre \acrshortpl{cu} y \acrshortpl{rf}.}
{l c c c c c c c c}
{matriz_trazabilidad}
{
	& \textbf{RF-1} & \textbf{RF-2} & \textbf{RF-3} & \textbf{RF-4} & \textbf{RF-5} & \textbf{RF-6} & \textbf{RF-7} & \textbf{RF-8} \\
}
{
	\textbf{CU-1} & X & & & & & & & \\
	\textbf{CU-2} & & X & X & X & X & & & \\
	\textbf{CU-3} & & & & & & X & & \\
	\textbf{CU-4} & & & & & & & X & \\
	\textbf{CU-5} & & & & & & & & X \\
}

